% Ollama相关书籍模板
% Ollama.tex

\documentclass[12pt,UTF8]{ctexbook}

% 设置纸张信息。
\input{../../../Book/common/page}

% 设置字体，并解决显示难检字问题。
\xeCJKsetup{AutoFallBack=true}
% 注意：ctexbook类已默认设置SimSun为CJKrmdefault
% 以下设置用于确保字体回退和扩展字体可用
% 仅设置扩展字体以避免与默认设置冲突
\setCJKfamilyfont{hei}{SimHei}
\setCJKfamilyfont{kai}{KaiTi}

% 目录 chapter 级别加点（.）。
\usepackage{titletoc}
\titlecontents{chapter}[0pt]{\vspace{3mm}\bf\addvspace{2pt}\filright}{\contentspush{\thecontentslabel\hspace{0.8em}}}{}{\titlerule*[8pt]{.}\contentspage}

% 支持目录点击跳转
\usepackage[colorlinks,linkcolor=blue,citecolor=blue,urlcolor=blue]{hyperref}

% 设置 part 和 chapter 标题格式。
\ctexset{
	part/name= {第,卷},
	part/number={\chinese{part}},
	chapter/name={第,篇},
	chapter/number={\arabic{chapter}}
}

% 图片相关设置。
\usepackage{graphicx}
\graphicspath{{Images/}}

% 设置署名格式。
\newenvironment{shuming}{\hfill\zihao{4}}

% 注脚每页重新编号，避免编号过大。
\usepackage[perpage]{footmisc}

% 列表项向右偏移。
\usepackage{enumitem}

\title{\heiti\zihao{0} Ollama}
\author{WangFei}
\date{\today}

\begin{document}

\maketitle
\tableofcontents

\frontmatter
\chapter{前言}

\mainmatter

\part{基础篇}
\chapter{Ollama简介}
\section{什么是Ollama}

Ollama 是使用大型语言模型（如 gpt-oss、Gemma 3、DeepSeek-R1、Qwen3 等）的最简单方法。

\section{Ollama的特点与优势}
\section{Ollama的应用场景}

\chapter{Ollama安装与配置}
\section{Windows系统安装}

在 PowerShell 中运行以下命令即可安装 Ollama：

\begin{verbatim}
irm https://ollama.com/install.ps1 | iex
\end{verbatim}

或者访问 https://ollama.com/download 下载 Ollama 安装程序。

\section{macOS系统安装}
\section{Linux系统安装}
\section{基本配置选项}

\chapter{Ollama命令行工具}
\section{基本命令使用}

Ollama 可在 macOS、Windows 和 Linux 上使用。在终端中运行 `ollama` 打开交互式菜单：

\begin{verbatim}
ollama
\end{verbatim}

使用 ↑/↓ 导航，按 enter 启动，→ 更改模型，esc 退出。

菜单提供快速访问：

\begin{itemize}
  \item Run a model - 开始交互式聊天
  \item Launch tools - Claude Code、Codex、OpenClaw 等
  \item Additional integrations - 在 "More…" 下可用
\end{itemize}

\section{模型管理}
\section{参数配置}

\part{实战篇}
\chapter{模型部署与使用}
\section{常用模型介绍}
\section{模型下载与管理}
\section{模型推理与对话}

使用 Ollama 启动开源模型进行编码工作，支持 Claude Code、Codex、OpenCode 等多种模型：

\begin{verbatim}
ollama launch claude
ollama launch codex
ollama launch opencode
\end{verbatim}

这将启动由最佳开源模型提供支持的编码环境。

\subsection{OpenClaw 开源 AI 助手}

OpenClaw 是一个具有 100+ 技能的个人 AI 助手，能够自动化工作、回答问题并处理任务，由开源模型提供支持：

\begin{verbatim}
ollama launch openclaw
\end{verbatim}

\chapter{Ollama API开发}
\section{REST API使用}

使用 API 将 Ollama 集成到您的应用程序中：
\begin{verbatim}
curl http://localhost:11434/api/chat -d '{ 
   "model": "gemma3", 
   "messages": [{ "role": "user", "content": "Hello!" }] 
 }'
\end{verbatim}

\subsection{Streaming 功能}
Streaming 允许您在模型生成文本时实时渲染文本。

\begin{itemize}
  \item REST API 默认启用 Streaming，但 SDK 默认禁用
  \item 在 SDK 中启用 Streaming，需要将 stream 参数设置为 True
\end{itemize}

\subsubsection{关键 Streaming 概念}
\begin{enumerate}
  \item \textbf{Chatting}：流式传输部分助手消息。每个块包含内容，因此您可以在消息到达时渲染它们。
  \item \textbf{Thinking}：支持思考的模型会在每个块中发出 thinking 字段。在流式块中检测此字段，以在最终答案到达前显示或隐藏推理轨迹。
  \item \textbf{Tool calling}：监视每个块中的 tool\_calls，执行请求的工具，并将工具输出附加回对话。
\end{enumerate}

\subsubsection{处理流式块}
需要累积部分字段以维护对话历史。这对于工具调用特别重要，因为模型的思考、工具调用和执行的工具结果必须在下次请求中传递回模型。

以下是处理流式块的 Python 代码示例：

\begin{verbatim}
from ollama import chat 

stream = chat( 
  model='qwen3', 
  messages=[{'role': 'user', 'content': 'What is 17 × 23?'}], 
  stream=True, 
)

in_thinking = False 
content = '' 
thinking = '' 
for chunk in stream: 
  if chunk.message.thinking: 
    if not in_thinking: 
      in_thinking = True 
      print('Thinking:\n', end='', flush=True) 
    print(chunk.message.thinking, end='', flush=True) 
    # accumulate the partial thinking 
    thinking += chunk.message.thinking 
  elif chunk.message.content: 
    if in_thinking: 
      in_thinking = False 
      print('\n\nAnswer:\n', end='', flush=True) 
    print(chunk.message.content, end='', flush=True) 
    # accumulate the partial content 
    content += chunk.message.content

  # append the accumulated fields to the messages for the next request 
  new_messages = [{ role: 'assistant', thinking: thinking, content: content }]
\end{verbatim}

以下是处理流式块的 JavaScript 代码示例：
\begin{verbatim}
import ollama from 'ollama' 

async function main() { 
  const stream = await ollama.chat({ 
    model: 'qwen3', 
    messages: [{ role: 'user', content: 'What is 17 × 23?' }], 
    stream: true, 
  }) 

  let inThinking = false 
  let content = '' 
  let thinking = '' 

  for await (const chunk of stream) { 
    if (chunk.message.thinking) { 
      if (!inThinking) { 
        inThinking = true 
        process.stdout.write('Thinking:\n') 
      } 
      process.stdout.write(chunk.message.thinking) 
      // accumulate the partial thinking 
      thinking += chunk.message.thinking 
    } else if (chunk.message.content) { 
      if (inThinking) { 
        inThinking = false 
        process.stdout.write('\n\nAnswer:\n') 
      } 
      process.stdout.write(chunk.message.content) 
      // accumulate the partial content 
      content += chunk.message.content 
    } 
  } 

  // append the accumulated fields to the messages for the next request 
  new_messages = [{ role: 'assistant', thinking: thinking, content: content }] 
} 

main().catch(console.error)
\end{verbatim}

\subsection{思考能力}
思考能力允许模型在生成最终答案之前展示其推理过程，这对于理解模型的思考过程非常有用。

\subsubsection{什么是思考能力}
支持思考的模型会发出一个思考字段，将其推理过程与最终答案分开。您可以使用此功能来审核模型步骤，在UI中为模型思考设置动画，或者在只需要最终响应时完全隐藏思考过程。

\subsubsection{支持的模型}
\begin{itemize}
  \item Qwen 3
  \item GPT-OSS（使用think级别：low、medium、high — 无法完全禁用思考过程）
  \item DeepSeek-v3.1
  \item DeepSeek R1
  \item 查看最新添加的思考模型
\end{itemize}

\subsubsection{在API调用中启用思考}
在聊天或生成请求中设置think字段。大多数模型接受布尔值（true/false）。而GPT-OSS则需要设置为low、medium或high来调整思考过程的长度。message.thinking（聊天端点）或thinking（生成端点）字段包含推理过程，而message.content / response包含最终答案。

\paragraph{cURL示例}
\begin{verbatim}
curl http://localhost:11434/api/chat -d '{
  "model": "qwen3",
  "messages": [{
    "role": "user",
    "content": "How many letter r are in strawberry?"
  }],
  "think": true,
  "stream": false
}'
\end{verbatim}

\paragraph{Python示例}
\begin{verbatim}
from ollama import chat

response = chat(
  model='qwen3',
  messages=[{'role': 'user', 'content': 'How many letter r are in strawberry?'}],
  think=True,
  stream=False,
)

print('Thinking:\n', response.message.thinking)
print('Answer:\n', response.message.content)
\end{verbatim}

\paragraph{JavaScript示例}
\begin{verbatim}
import ollama from 'ollama'

const response = await ollama.chat({
  model: 'deepseek-r1',
  messages: [{ role: 'user', content: 'How many letter r are in strawberry?' }],
  think: true,
  stream: false,
})

console.log('Thinking:\n', response.message.thinking)
console.log('Answer:\n', response.message.content)
\end{verbatim}

注意：GPT-OSS需要将think设置为"low"、"medium"或"high"。对于该模型，传递true/false会被忽略。

\subsubsection{处理思考过程}
在流式响应中，您可以通过检查chunk.message.thinking字段来识别思考过程，类似于之前的流式处理示例。

\subsubsection{流式处理思考过程}
思考流会在答案令牌之前交错思考令牌。检测第一个思考块以渲染"思考"部分，然后在message.content到达时切换到最终回复。

\paragraph{Python示例}
\begin{verbatim}
from ollama import chat

stream = chat(
  model='qwen3',
  messages=[{'role': 'user', 'content': 'What is 17 × 23?'}],
  think=True,
  stream=True,
)

in_thinking = False

for chunk in stream:
  if chunk.message.thinking and not in_thinking:
    in_thinking = True
    print('Thinking:\n', end='')

  if chunk.message.thinking:
    print(chunk.message.thinking, end='')
  elif chunk.message.content:
    if in_thinking:
      print('\n\nAnswer:\n', end='')
      in_thinking = False
    print(chunk.message.content, end='')
\end{verbatim}

\paragraph{JavaScript示例}
\begin{verbatim}
import ollama from 'ollama'

async function main() {
  const stream = await ollama.chat({
    model: 'qwen3',
    messages: [{ role: 'user', content: 'What is 17 × 23?' }],
    think: true,
    stream: true,
  })

  let inThinking = false

  for await (const chunk of stream) {
    if (chunk.message.thinking && !inThinking) {
      inThinking = true
      process.stdout.write('Thinking:\n')
    }

    if (chunk.message.thinking) {
      process.stdout.write(chunk.message.thinking)
    } else if (chunk.message.content) {
      if (inThinking) {
        process.stdout.write('\n\nAnswer:\n')
        inThinking = false
      }
      process.stdout.write(chunk.message.content)
    }
  }
}

main()
\end{verbatim}

\subsubsection{CLI快速参考}
在命令行中使用思考能力的快速参考：

\begin{verbatim}
# 为单次运行启用思考能力
ollama run deepseek-r1 --think "Where should I visit in Lisbon?"

# 禁用思考能力
ollama run deepseek-r1 --think=false "Summarize this article"

# 隐藏思考过程但仍然使用思考模型
ollama run deepseek-r1 --hidethinking "Is 9.9 bigger or 9.11?"

# 在交互式会话中，使用 /set think 或 /set nothink 切换

# GPT-OSS 只接受级别设置
ollama run gpt-oss --think=low "Draft a headline"  # 可将 low 替换为 medium 或 high
\end{verbatim}

注意：对于支持思考能力的模型，思考功能在 CLI 和 API 中默认启用。

\section{结构化输出}
结构化输出允许您在模型响应上强制执行 JSON 模式，以便您可以可靠地提取结构化数据、描述图像或保持每个回复的一致性。

\subsection{生成结构化JSON}

\paragraph{cURL示例}
\begin{verbatim}
curl -X POST http://localhost:11434/api/chat -H "Content-Type: application/json" -d '{
  "model": "gpt-oss",
  "messages": [{"role": "user", "content": "Tell me about Canada in one line"}],
  "stream": false,
  "format": "json"
}'
\end{verbatim}

\paragraph{Python示例}
\begin{verbatim}
from ollama import chat

response = chat(
  model='gpt-oss',
  messages=[{'role': 'user', 'content': 'Tell me about Canada.'}],
  format='json'
)

print(response.message.content)
\end{verbatim}

\paragraph{JavaScript示例}
\begin{verbatim}
import ollama from 'ollama'

const response = await ollama.chat({
  model: 'gpt-oss',
  messages: [{ role: 'user', content: 'Tell me about Canada.' }],
  format: 'json'
})

console.log(response.message.content)
\end{verbatim}

\subsection{使用架构生成结构化JSON}
您可以在format字段中提供JSON架构。理想情况下，您还应该在提示中作为字符串传递JSON架构，以确保模型响应的准确性。

\paragraph{JavaScript示例（使用Zod架构）}
使用zodToJsonSchema()序列化Zod架构并解析结构化响应：
\begin{verbatim}
import ollama from 'ollama'
import { z } from 'zod'
import { zodToJsonSchema } from 'zod-to-json-schema'

const Country = z.object({
  name: z.string(),
  capital: z.string(),
  languages: z.array(z.string()),
})

const response = await ollama.chat({
  model: 'gpt-oss',
  messages: [{ role: 'user', content: 'Tell me about Canada.' }],
  format: zodToJsonSchema(Country),
})

const country = Country.parse(JSON.parse(response.message.content))
console.log(country)
\end{verbatim}

\paragraph{Python示例（使用Pydantic）}
Use Pydantic models and pass model_json_schema() to format, then validate the response:
\begin{verbatim}
from ollama import chat
from pydantic import BaseModel

class Country(BaseModel):
    name: str
    capital: str
    languages: list[str]

response = chat(
    model='gpt-oss',
    messages=[{"role": "user", "content": "Tell me about Canada."}],
    format=Country.model_json_schema(),
)

country = Country.model_validate_json(response.message.content)
print(country)
\end{verbatim}

\paragraph{cURL示例（使用JSON架构）}
\begin{verbatim}
curl -X POST http://localhost:11434/api/chat -H "Content-Type: application/json" -d '{
  "model": "gpt-oss",
  "messages": [{"role": "user", "content": "Tell me about Canada."}],
  "stream": false,
  "format": {
    "type": "object",
    "properties": {
      "name": {"type": "string"},
      "capital": {"type": "string"},
      "languages": {
        "type": "array",
        "items": {"type": "string"}
      }
    },
    "required": ["name", "capital", "languages"]
  }
}'
\end{verbatim}

\subsection{结构化输出示例}

\subsubsection{示例：提取结构化数据}
定义您想要返回的对象，让模型填充字段：

\begin{verbatim}
from ollama import chat
from pydantic import BaseModel

class Pet(BaseModel):
  name: str
  animal: str
  age: int
  color: str | None
  favorite_toy: str | None

class PetList(BaseModel):
  pets: list[Pet]

response = chat(
  model='gpt-oss',
  messages=[{'role': 'user', 'content': 'I have two cats named Luna and Loki...'}],
  format=PetList.model_json_schema(),
)

pets = PetList.model_validate_json(response.message.content)
print(pets)
\end{verbatim}

\subsubsection{示例：带结构化输出的视觉模型}
视觉模型接受相同的format参数，实现对图像的确定性描述：

\begin{verbatim}
from ollama import chat
from pydantic import BaseModel
from typing import Literal, Optional

class Object(BaseModel):
  name: str
  confidence: float
  attributes: str

class ImageDescription(BaseModel):
  summary: str
  objects: list[Object]
  scene: str
  colors: list[str]
  time_of_day: Literal['Morning', 'Afternoon', 'Evening', 'Night']
  setting: Literal['Indoor', 'Outdoor', 'Unknown']
  text_content: Optional[str] = None

response = chat(
  model='gemma3',
  messages=[{
    'role': 'user',
    'content': 'Describe this photo and list the objects you detect.',
    'images': ['path/to/image.jpg'],
  }],
  format=ImageDescription.model_json_schema(),
  options={'temperature': 0},
)

image_description = ImageDescription.model_validate_json(response.message.content)
print(image_description)
\end{verbatim}

\subsubsection{Tips for reliable structured outputs}

\begin{itemize}
  \item Define schemas with Pydantic (Python) or Zod (JavaScript) so they can be reused for validation
  \item Lower the temperature (e.g., set it to 0) for more deterministic completions
  \item Structured outputs work through the OpenAI-compatible API via response_format
\end{itemize}

\section{视觉模型}
视觉模型接受图像和文本，以便模型可以描述、分类和回答关于它所看到的内容的问题。

\subsection{快速开始}

\begin{verbatim}
ollama run gemma3 ./image.png whats in this image?
\end{verbatim}

\subsection{使用Ollama的API}
提供一个images数组。SDK接受文件路径、URL或原始字节，而REST API需要base64编码的图像数据。

\subsubsection{cURL示例}
\begin{verbatim}
# 1. 下载示例图像
curl -L -o test.jpg "https://upload.wikimedia.org/wikipedia/commons/3/3a/Cat03.jpg"

# 2. 编码图像
IMG=$(base64 < test.jpg | tr -d '\n')

# 3. 发送到Ollama
curl -X POST http://localhost:11434/api/chat \
-H "Content-Type: application/json" \
-d '{
    "model": "gemma3",
    "messages": [{
    "role": "user",
    "content": "What is in this image?",
    "images": ["'"$IMG"'"]
    }],
    "stream": false
}'
\end{verbatim}

\subsubsection{Python示例}
\begin{verbatim}
from ollama import chat
# from pathlib import Path

# Pass in the path to the image
path = input('Please enter the path to the image: ')

# You can also pass in base64 encoded image data
# img = base64.b64encode(Path(path).read_bytes()).decode()
# or the raw bytes
# img = Path(path).read_bytes()

response = chat(
  model='gemma3',
  messages=[
    {
      'role': 'user',
      'content': 'What is in this image? Be concise.',
      'images': [path],
    }
  ],
)

print(response.message.content)
\end{verbatim}

\subsubsection{JavaScript示例}
\begin{verbatim}
import ollama from 'ollama'

const imagePath = '/absolute/path/to/image.jpg'
const response = await ollama.chat({
  model: 'gemma3',
  messages: [
    { role: 'user', content: 'What is in this image?', images: [imagePath] }
  ],
  stream: false,
})

console.log(response.message.content)
\end{verbatim}

\section{Embeddings}

生成文本embeddings用于语义搜索、检索和RAG。

Embeddings将文本转换为数值向量，您可以将其存储在向量数据库中，使用余弦相似度进行搜索，或在RAG管道中使用。向量长度取决于模型（通常为384-1024维）。

\subsection{推荐模型}

\begin{itemize}
  \item embeddinggemma
  \item qwen3-embedding
  \item all-minilm
\end{itemize}

\subsection{生成Embeddings}

\subsubsection{CLI示例}
\begin{verbatim}
# 直接从命令行生成embeddings
ollama run embeddinggemma "Hello world"

# 也可以通过管道传递文本生成embeddings
echo "Hello world" | ollama run embeddinggemma
\end{verbatim}

输出是一个JSON数组。/api/embed端点返回L2归一化（单位长度）的向量。

\subsubsection{cURL示例}
\begin{verbatim}
curl -X POST http://localhost:11434/api/embed \
  -H "Content-Type: application/json" \
  -d '{
    "model": "embeddinggemma",
    "input": "The quick brown fox jumps over the lazy dog."
  }'
\end{verbatim}

\subsubsection{Python示例}
\begin{verbatim}
import ollama

single = ollama.embed(
  model='embeddinggemma',
  input='The quick brown fox jumps over the lazy dog.'
)
print(len(single['embeddings'][0]))  # vector length
\end{verbatim}

\subsubsection{JavaScript示例}
\begin{verbatim}
import ollama from 'ollama'

const single = await ollama.embed({
  model: 'embeddinggemma',
  input: 'The quick brown fox jumps over the lazy dog.',
})
console.log(single.embeddings[0].length) // vector length
\end{verbatim}

\subsubsection{Generate a batch of embeddings}
Pass an array of strings to input:

\paragraph{cURL example}
\begin{verbatim}
curl -X POST http://localhost:11434/api/embed \
  -H "Content-Type: application/json" \
  -d '{
    "model": "embeddinggemma",
    "input": [
      "The quick brown fox jumps over the lazy dog.",
      "The five boxing wizards jump quickly.",
      "Jackdaws love my big sphinx of quartz."
    ]
  }'
\end{verbatim}

\paragraph{Python example}
\begin{verbatim}
import ollama

batch = ollama.embed(
  model='embeddinggemma',
  input=[
    'The quick brown fox jumps over the lazy dog.',
    'The five boxing wizards jump quickly.',
    'Jackdaws love my big sphinx of quartz.',
  ],
)
print(len(batch['embeddings']))  # number of vectors
\end{verbatim}

\paragraph{JavaScript example}
\begin{verbatim}
import ollama from 'ollama'

const batch = await ollama.embed({
  model: 'embeddinggemma',
  input: [
    'The quick brown fox jumps over the lazy dog.',
    'The five boxing wizards jump quickly.',
    'Jackdaws love my big sphinx of quartz.',
  ],
})
console.log(batch.embeddings.length) // number of vectors
\end{verbatim}

\subsubsection{Tips}

\begin{itemize}
  \item Use cosine similarity for most semantic search use cases.
  \item Use the same embedding model for both indexing and querying.
\end{itemize}

\section{Tool calling}

Ollama supports tool calling (also known as function calling) which allows a model to invoke tools and incorporate their results into its replies.

\subsection{Calling a single tool}
Invoke a single tool and include its response in a follow-up request. Also known as "single-shot" tool calling.

\subsubsection{cURL example}
\begin{verbatim}
curl -s http://localhost:11434/api/chat -H "Content-Type: application/json" -d '{
  "model": "qwen3",
  "messages": [{"role": "user", "content": "What is the temperature in New York?"}],
  "stream": false,
  "tools": [
    {
      "type": "function",
      "function": {
        "name": "get_temperature",
        "description": "Get the current temperature for a city",
        "parameters": {
          "type": "object",
          "required": ["city"],
          "properties": {
            "city": {"type": "string", "description": "The name of the city"}
          }
        }
      }
    }
  ]
}'
\end{verbatim}

\subsubsection{Generate a response with a single tool result}
\begin{verbatim}
curl -s http://localhost:11434/api/chat -H "Content-Type: application/json" -d '{
  "model": "qwen3",
  "messages": [
    {"role": "user", "content": "What is the temperature in New York?"},
    {
      "role": "assistant",
      "tool_calls": [
        {
          "type": "function",
          "function": {
            "index": 0,
            "name": "get_temperature",
            "arguments": {"city": "New York"}
          }
        }
      ]
    },
    {"role": "tool", "tool_name": "get_temperature", "content": "22°C"}
  ],
  "stream": false
}'
\end{verbatim}

\subsubsection{Python example}

Install the Ollama Python SDK:

\begin{verbatim}
# with pip
pip install ollama -U

# with uv
uv add ollama
\end{verbatim}

\begin{verbatim}
from ollama import chat

def get_temperature(city: str) -> str:
  """Get the current temperature for a city
  
  Args:
    city: The name of the city

  Returns:
    The current temperature for the city
  """
  temperatures = {
    "New York": "22°C",
    "London": "15°C",
    "Tokyo": "18°C",
  }
  return temperatures.get(city, "Unknown")

messages = [{"role": "user", "content": "What is the temperature in New York?"}]

# pass functions directly as tools in the tools list or as a JSON schema
response = chat(model="qwen3", messages=messages, tools=[get_temperature], think=True)

messages.append(response.message)
if response.message.tool_calls:
  # only recommended for models which only return a single tool call
  call = response.message.tool_calls[0]
  result = get_temperature(**call.function.arguments)
  # add the tool result to the messages
  messages.append({"role": "tool", "tool_name": call.function.name, "content": str(result)})

  final_response = chat(model="qwen3", messages=messages, tools=[get_temperature], think=True)
  print(final_response.message.content)
\end{verbatim}

\subsubsection{JavaScript example}

Install the Ollama JavaScript library:

\begin{verbatim}
# with npm
npm i ollama

# with bun
bun i ollama
\end{verbatim}

\begin{verbatim}
import ollama from 'ollama'

function getTemperature(city: string): string {
  const temperatures: Record<string, string> = {
    'New York': '22°C',
    'London': '15°C',
    'Tokyo': '18°C',
  }
  return temperatures[city] ?? 'Unknown'
}

const tools = [
  {
    type: 'function',
    function: {
      name: 'get_temperature',
      description: 'Get the current temperature for a city',
      parameters: {
        type: 'object',
        required: ['city'],
        properties: {
          city: { type: 'string', description: 'The name of the city' },
        },
      },
    },
  },
]

const messages = [{ role: 'user', content: "What is the temperature in New York?" }]

const response = await ollama.chat({
  model: 'qwen3',
  messages,
  tools,
  think: true,
})

messages.push(response.message)
if (response.message.tool_calls?.length) {
  // only recommended for models which only return a single tool call
  const call = response.message.tool_calls[0]
  const args = call.function.arguments as { city: string }
  const result = getTemperature(args.city)
  // add the tool result to the messages
  messages.push({ role: 'tool', tool_name: call.function.name, content: result })

  // generate the final response
  const finalResponse = await ollama.chat({ model: 'qwen3', messages, tools, think: true })
  console.log(finalResponse.message.content)
}
\end{verbatim}

\subsection{Parallel tool calling}

Request multiple tool calls in parallel, then send all tool responses back to the model.

\subsubsection{cURL example}
\begin{verbatim}
curl -s http://localhost:11434/api/chat -H "Content-Type: application/json" -d '{
  "model": "qwen3",
  "messages": [{"role": "user", "content": "What are the current weather conditions and temperature in New York and London?"}],
  "stream": false,
  "tools": [
    {
      "type": "function",
      "function": {
        "name": "get_temperature",
        "description": "Get the current temperature for a city",
        "parameters": {
          "type": "object",
          "required": ["city"],
          "properties": {
            "city": {"type": "string", "description": "The name of the city"}
          }
        }
      }
    },
    {
      "type": "function",
      "function": {
        "name": "get_conditions",
        "description": "Get the current weather conditions for a city",
        "parameters": {
          "type": "object",
          "required": ["city"],
          "properties": {
            "city": {"type": "string", "description": "The name of the city"}
          }
        }
      }
    }
  ]
}'
\end{verbatim}

\subsubsection{Generate a response with multiple tool results}
\begin{verbatim}
curl -s http://localhost:11434/api/chat -H "Content-Type: application/json" -d '{
  "model": "qwen3",
  "messages": [
    {"role": "user", "content": "What are the current weather conditions and temperature in New York and London?"},
    {
      "role": "assistant",
      "tool_calls": [
        {
          "type": "function",
          "function": {
            "index": 0,
            "name": "get_temperature",
            "arguments": {"city": "New York"}
          }
        },
        {
          "type": "function",
          "function": {
            "index": 1,
            "name": "get_conditions",
            "arguments": {"city": "New York"}
          }
        },
        {
          "type": "function",
          "function": {
            "index": 2,
            "name": "get_temperature",
            "arguments": {"city": "London"}
          }
        },
        {
          "type": "function",
          "function": {
            "index": 3,
            "name": "get_conditions",
            "arguments": {"city": "London"}
          }
        }
      ]
    },
    {"role": "tool", "tool_name": "get_temperature", "content": "22°C"},
    {"role": "tool", "tool_name": "get_conditions", "content": "Partly cloudy"},
    {"role": "tool", "tool_name": "get_temperature", "content": "15°C"},
    {"role": "tool", "tool_name": "get_conditions", "content": "Rainy"}
  ],
  "stream": false
}'
\end{verbatim}

\subsubsection{Python example}
\begin{verbatim}
from ollama import chat

def get_temperature(city: str) -> str:
  """Get the current temperature for a city
  
  Args:
    city: The name of the city

  Returns:
    The current temperature for the city
  """
  temperatures = {
    "New York": "22°C",
    "London": "15°C",
    "Tokyo": "18°C"
  }
  return temperatures.get(city, "Unknown")

def get_conditions(city: str) -> str:
  """Get the current weather conditions for a city
  
  Args:
    city: The name of the city

  Returns:
    The current weather conditions for the city
  """
  conditions = {
    "New York": "Partly cloudy",
    "London": "Rainy",
    "Tokyo": "Sunny"
  }
  return conditions.get(city, "Unknown")


messages = [{'role': 'user', 'content': 'What are the current weather conditions and temperature in New York and London?'}]

# The python client automatically parses functions as a tool schema so we can pass them directly
# Schemas can be passed directly in the tools list as well
response = chat(model='qwen3', messages=messages, tools=[get_temperature, get_conditions], think=True)

# add the assistant message to the messages
messages.append(response.message)
if response.message.tool_calls:
  # process each tool call
  for call in response.message.tool_calls:
    # execute the appropriate tool
    if call.function.name == 'get_temperature':
      result = get_temperature(**call.function.arguments)
    elif call.function.name == 'get_conditions':
      result = get_conditions(**call.function.arguments)
    else:
      result = 'Unknown tool'
    # add the tool result to the messages
    messages.append({'role': 'tool',  'tool_name': call.function.name, 'content': str(result)})

  # generate the final response
  final_response = chat(model='qwen3', messages=messages, tools=[get_temperature, get_conditions], think=True)
  print(final_response.message.content)
\end{verbatim}

\subsubsection{JavaScript example}

Install the Ollama JavaScript library:

\begin{verbatim}
# with npm
npm i ollama

# with bun
bun i ollama
\end{verbatim}

\begin{verbatim}
import ollama from 'ollama'

function getTemperature(city: string): string {
  const temperatures: Record<string, string> = {
    'New York': '22°C',
    'London': '15°C',
    'Tokyo': '18°C',
  }
  return temperatures[city] ?? 'Unknown'
}

const tools = [
  {
    type: 'function',
    function: {
      name: 'get_temperature',
      description: 'Get the current temperature for a city',
      parameters: {
        type: 'object',
        required: ['city'],
        properties: {
          city: { type: 'string', description: 'The name of the city' },
        },
      },
    },
  },
]

const messages = [{ role: 'user', content: "What is the temperature in New York?" }]

const response = await ollama.chat({
  model: 'qwen3',
  messages,
  tools,
  think: true,
})

messages.push(response.message)
if (response.message.tool_calls?.length) {
  // only recommended for models which only return a single tool call
  const call = response.message.tool_calls[0]
  const args = call.function.arguments as { city: string }
  const result = getTemperature(args.city)
  // add the tool result to the messages
  messages.push({ role: 'tool', tool_name: call.function.name, content: result })

  // generate the final response
  const finalResponse = await ollama.chat({ model: 'qwen3', messages, tools, think: true })
  console.log(finalResponse.message.content)
}
\end{verbatim}

\section{CLI Reference}

\subsection{Run a model}
\begin{verbatim}
ollama run gemma3
\end{verbatim}

\subsection{Launch integrations}
\begin{verbatim}
ollama launch
\end{verbatim}

Configure and launch external applications to use Ollama models. This provides an interactive way to set up and start integrations with supported apps.

\subsubsection{Supported integrations}
\begin{itemize}
  \item OpenCode - Open-source coding assistant
  \item Claude Code - Anthropic's agentic coding tool
  \item Codex - OpenAI's coding assistant
  \item Droid - Factory's AI coding agent
\end{itemize}

\subsubsection{Examples}
\begin{itemize}
  \item Launch an integration interactively: 
    \begin{verbatim}
    ollama launch
    \end{verbatim}
  \item Launch a specific integration: 
    \begin{verbatim}
    ollama launch claude
    \end{verbatim}
  \item Launch with a specific model: 
    \begin{verbatim}
    ollama launch claude --model qwen3-coder
    \end{verbatim}
  \item Configure without launching: 
    \begin{verbatim}
    ollama launch droid --config
    \end{verbatim}
\end{itemize}

\subsection{Multiline input}
For multiline input, you can wrap text with """:
\begin{verbatim}
>>> """Hello,
... world!
... """
I'm a basic program that prints the famous "Hello, world!" message to the console.
\end{verbatim}

\subsection{Multimodal models}
\begin{verbatim}
ollama run gemma3 "What's in this image? /Users/jmorgan/Desktop/smile.png"
\end{verbatim}

\subsection{Generate embeddings}
\begin{verbatim}
ollama run embeddinggemma "Hello world"
\end{verbatim}

Output is a JSON array:
\begin{verbatim}
echo "Hello world" | ollama run nomic-embed-text
\end{verbatim}

\subsection{Download a model}
\begin{verbatim}
ollama pull gemma3
\end{verbatim}

\subsection{Remove a model}
\begin{verbatim}
ollama rm gemma3
\end{verbatim}

\subsection{List models}
\begin{verbatim}
ollama ls
\end{verbatim}

\subsection{Sign in to Ollama}
\begin{verbatim}
ollama signin
\end{verbatim}

\subsection{Sign out of Ollama}
\begin{verbatim}
ollama signout
\end{verbatim}

\subsection{Create a customized model}
First, create a Modelfile:
\begin{verbatim}
FROM gemma3
SYSTEM """You are a happy cat."""
\end{verbatim}

Then run ollama create:
\begin{verbatim}
ollama create -f Modelfile
\end{verbatim}

\subsection{List running models}
\begin{verbatim}
ollama ps
\end{verbatim}

\subsection{Stop a running model}
\begin{verbatim}
ollama stop gemma3
\end{verbatim}

\subsection{Start Ollama}
\begin{verbatim}
ollama serve
\end{verbatim}

To view a list of environment variables that can be set run ollama serve --help

\section{API参考}

\subsection{简介}

Ollama的API允许您以编程方式运行和与模型交互。

\subsection{快速开始}

如果您是刚开始使用，请按照快速入门文档来启动和运行Ollama的API。

\subsection{基础URL}

安装后，Ollama的API默认在以下地址提供服务：

\begin{verbatim}
http://localhost:11434/api
\end{verbatim}

对于在ollama.com上运行的云模型，相同的API可通过以下基础URL访问：

\begin{verbatim}
https://ollama.com/api
\end{verbatim}

\subsection{示例请求}

Ollama运行后，其API会自动可用，可通过curl访问：

\begin{verbatim}
curl http://localhost:11434/api/generate -d '{
  "model": "gemma3",
  "prompt": "为什么天空是蓝色的？"
}'
\end{verbatim}

\subsection{库}

Ollama为Python和JavaScript提供官方库：

\begin{itemize}
  \item Python
  \item JavaScript
\end{itemize}

Ollama还有几个社区维护的库。完整列表请参见Ollama GitHub仓库。

\subsection{版本控制}

Ollama的API没有严格的版本控制，但API预计是稳定且向后兼容的。弃用很少见，将在发布说明中宣布。

\subsection{认证}

通过 http://localhost:11434 本地访问Ollama的API时不需要认证。以下情况需要认证：

\begin{itemize}
  \item 通过ollama.com运行云模型
  \item 发布模型
  \item 下载私有模型
\end{itemize}

Ollama支持两种认证方法：

\begin{itemize}
  \item 登录：从本地安装登录，Ollama将在运行命令时自动处理对ollama.com的请求认证
  \item API密钥：用于以编程方式访问ollama.com的API
\end{itemize}

\subsubsection{登录}

要从本地Ollama安装登录到ollama.com，请运行：

\begin{verbatim}
ollama signin
\end{verbatim}

登录后，Ollama将根据需要自动认证命令：

\begin{verbatim}
ollama run gpt-oss:120b-cloud
\end{verbatim}

同样，当访问需要云访问的本地API端点时，Ollama将自动认证请求：

\begin{verbatim}
curl http://localhost:11434/api/generate -d '{
  "model": "gpt-oss:120b-cloud",
  "prompt": "为什么天空是蓝色的？"
}'
\end{verbatim}

\subsubsection{API密钥}

对于直接访问 `https://ollama.com/api` 提供的ollama.com API，需要通过API密钥进行认证。首先，创建一个API密钥，然后设置OLLAMA_API_KEY环境变量：

\begin{verbatim}
export OLLAMA_API_KEY=your_api_key
\end{verbatim}

然后在Authorization头中使用API密钥：

\begin{verbatim}
curl https://ollama.com/api/generate \
  -H "Authorization: Bearer $OLLAMA_API_KEY" \
  -d '{
    "model": "gpt-oss:120b",
    "prompt": "为什么天空是蓝色的？",
    "stream": false
  }'
\end{verbatim}

API密钥目前不会过期，但您可以随时在API密钥设置中撤销它们。

\subsection{流式传输}

某些API端点默认会流式传输响应，例如 `/api/generate`。这些响应以换行分隔的JSON格式（即 `application/x-ndjson` 内容类型）提供。例如：

\begin{verbatim}
{"model":"gemma3","created_at":"2025-10-26T17:15:24.097767Z","response":"That","done":false}
{"model":"gemma3","created_at":"2025-10-26T17:15:24.109172Z","response":"'","done":false}
{"model":"gemma3","created_at":"2025-10-26T17:15:24.121485Z","response":"s","done":false}
{"model":"gemma3","created_at":"2025-10-26T17:15:24.132802Z","response":" a","done":false}
{"model":"gemma3","created_at":"2025-10-26T17:15:24.143931Z","response":" fantastic","done":false}
{"model":"gemma3","created_at":"2025-10-26T17:15:24.155176Z","response":" question","done":false}
{"model":"gemma3","created_at":"2025-10-26T17:15:24.166576Z","response":"!","done":true, "done_reason": "stop"}
\end{verbatim}

\subsubsection{禁用流式传输}

可以通过在支持流式传输的任何端点的请求体中提供 `{"stream": false}` 来禁用流式传输。这将导致响应以 `application/json` 格式返回：

\begin{verbatim}
{"model":"gemma3","created_at":"2025-10-26T17:15:24.166576Z","response":"That's a fantastic question!","done":true}
\end{verbatim}

\subsubsection{何时使用流式传输与非流式传输}

流式传输（默认）：

\begin{itemize}
  \item 实时响应生成
  \item 更低的感知延迟
  \item 更适合长生成
\end{itemize}

非流式传输：

\begin{itemize}
  \item 处理更简单
  \item 更适合短响应或结构化输出
  \item 在某些应用中更容易处理
\end{itemize}

\subsection{使用指标}

Ollama的API响应包含可用于衡量性能和模型使用情况的指标：

\begin{itemize}
  \item total_duration: 响应生成所需的总时间
  \item load_duration: 模型加载所需的时间
  \item prompt_eval_count: 处理的输入令牌数量
  \item prompt_eval_duration: 评估提示所需的时间
  \item eval_count: 处理的输出令牌数量
  \item eval_duration: 生成输出令牌所需的时间
\end{itemize}

所有时间值均以纳秒为单位。

\subsubsection{示例响应}

对于返回使用指标的端点，响应体将包含使用字段。例如，对 `/api/generate` 的非流式调用可能返回以下响应：

\begin{verbatim}
{
  "model": "gemma3",
  "created_at": "2025-10-17T23:14:07.414671Z",
  "response": "Hello! How can I help you today?",
  "done": true,
  "done_reason": "stop",
  "total_duration": 174560334,
  "load_duration": 101397084,
  "prompt_eval_count": 11,
  "prompt_eval_duration": 13074791,
  "eval_count": 18,
  "eval_duration": 52479709
}
\end{verbatim}

对于返回流式响应的端点，使用字段包含在最后一个块中，其中 `done` 为 `true`。

\section{SDK集成}
\section{示例代码}

\chapter{Ollama与其他工具集成}
\section{与LangChain集成}
\section{与LLM框架集成}
\section{与应用系统集成}

\subsection{运行任何应用或代理}

使用 Ollama，您可以将最新的开源模型连接到您喜欢的应用程序或代理，轻松在它们之间切换：

\begin{verbatim}
ollama
\end{verbatim}

这使得您可以灵活地为不同的任务选择最适合的模型，而无需修改应用程序代码。

\part{高级篇}
\chapter{Ollama架构与原理}
\section{系统架构}
\section{模型加载机制}
\section{推理优化技术}

\chapter{自定义模型开发}
\section{模型微调}
\section{模型量化}
\section{模型部署最佳实践}

\chapter{Ollama性能优化}
\section{硬件资源配置}
\section{推理速度优化}
\section{内存使用优化}

\backmatter

\end{document}